% year: תשע"ט
% semester: א
% moed: א
% number: 3א
% points: 10
% categories: func-limits, proofs
% source: exams/מבחנים/תשעט/מועד א תשעט.pdf

\begin{question}
תהיה פונקציה $f(x)$ מוגדרת בסביבת נקודה $x_0$. השלימו את ההגדרה הבאה:

\textbf{מספר $a$ הוא גבול של $f(x)$ כאשר $x$ שואף ל-$x_0$ אם לכל \ldots}
\end{question}

\begin{solution}
(הגדרת קושי, $\eps$--$\delta$) מספר $a$ הוא גבול של $f(x)$ כאשר $x$ שואף ל-$x_0$ אם לכל $\eps>0$ קיים $\delta>0$ כך שלכל $x$ המקיים $0<|x-x_0|<\delta$ מתקיים $|f(x)-a|<\eps$.

(הגדרה שקולה לפי היינה: לכל סדרה $x_n\to x_0$ עם $x_n\neq x_0$ מתקיים $f(x_n)\to a$.)
\end{solution}
